\documentclass[11pt]{article}
\usepackage[utf8]{inputenc}
\usepackage{amsmath, amssymb}
\usepackage{geometry}
\usepackage{listings}
\usepackage{xcolor}
\usepackage{graphicx}
\usepackage{tabularx}
\usepackage{booktabs}
\usepackage{hyperref}

\geometry{a4paper, margin=1in}

\definecolor{codegray}{rgb}{0.5,0.5,0.5}
\definecolor{codepurple}{rgb}{0.58,0,0.82}
\definecolor{backcolour}{rgb}{0.95,0.95,0.92}

\lstset{
    language=Python,
    backgroundcolor=\color{backcolour},
    commentstyle=\color{codegray},
    keywordstyle=\color{magenta},
    stringstyle=\color{codepurple},
    basicstyle=\ttfamily\small,
    breaklines=true,
    showstringspaces=false
}

\title{Math Lab: Abstract Algebra and QR Codes \\ \large Problem Set \& Implementation Guide}
\author{VB}
\date{\today}

\begin{document}

\maketitle

\section*{Problem 0: Start the engine}

Consider the files given.
\begin{center}
\renewcommand{\arraystretch}{1.5}
\begin{tabularx}{\textwidth}{@{} >{\bfseries\raggedright\arraybackslash}p{5.5cm} X @{}}
\toprule
\textcolor{blue}{Filename} & \textcolor{blue}{Purpose} \\ \midrule
\multicolumn{2}{@{}l}{\textbf{Direct Interaction}} \\ \midrule
\lstinline|implementation_tasks.py| & \textcolor{red}{Here you write your code.} Contains the core abstract algebra (Galois Fields) and Reed-Solomon coding logic intended for student implementation. \\
\lstinline|qr_lab_dashboard.py| & \textcolor{red}{This file you run.} A dedicated educational sandbox for debugging individual math levels through an interactive UI. \\
\lstinline|qr_code_demo.py| & \textbf{May be run standalone after you solve levels 1-9.} The main engine that ties all your math together to generate a fully functioning, scannable physical QR code. \\
\lstinline|levels/level_*.py| & Unittests and graphical representation for each specific problem. \textbf{Can be run standalone.} Use \lstinline|python levels/level_*.py| for the graphical debugger or \lstinline|python levels/level_*.py --no-graphics| for unittests only. \\ \midrule
\multicolumn{2}{@{}l}{\textbf{Under the Hood}} \\ \midrule
\lstinline|levels/base_level.py| & Base class for graphical debuggers. \\
\lstinline|qr_code.py| & Pre-baked ISO formatting logic (BCH codes, masking) that utilizes your math. \\
\bottomrule
\end{tabularx}
\end{center}

Run the file \lstinline|qr_lab_dashboard.py|. You should see the interactive dashboard.
Now use any editor to open the file \lstinline|implementation_tasks.py|. You should see the API for the functions you must implement.

\begin{lstlisting}
# =============================================================
#  STUDENT IMPLEMENTATION
# =============================================================

def extended_gcd(a, b):
    pass

def mod_inverse(a, m):
    pass

<<<< MORE CODE >>>>
\end{lstlisting}

First solve the math part so you understand the underlying finite fields. Then implement the code. After implementation, click "Run All Tests" in the dashboard. If a light turns green, your math is correct. Click the level buttons to see a visual derivation of your logic.


\section*{Problem 1: Prime Fields GF(p)}

\subsubsection*{Mathematical part}
The field $GF(p)$ is obtained by performing arithmetic modulo a prime number. A fundamental theorem states that $\mathbb Z/p\mathbb Z$ is a field if and only if $p$ is prime. This guarantees that every nonzero element has a unique inverse, and polynomial long division works exactly as over the real numbers.

\textbf{Worked Example: Polynomial Division in $GF(3)$} \\
Let us divide $f(x) = 2x^2 + 1$ by $g(x) = x + 2$ over $GF(3)$.
\begin{enumerate}
    \item \textbf{Leading terms:} Divide $2x^2$ by $x$ to get $2x$.
    \item \textbf{Multiply \& Subtract:} $2x \cdot (x + 2) = 2x^2 + 4x$. Modulo 3, this is $2x^2 + x$.
    \item $(2x^2 + 1) - (2x^2 + x) = -x + 1$. Modulo 3, $-x$ becomes $2x$. The remainder is $2x + 1$.
    \item \textbf{Repeat:} Divide $2x$ by $x$ to get $2$.
    \item $2 \cdot (x + 2) = 2x + 4 \equiv 2x + 1 \pmod 3$.
    \item $(2x + 1) - (2x + 1) = 0$.
\end{enumerate}
\textbf{Result:} $q(x) = 2x + 2$, $r(x) = 0$.

\subsubsection*{Implementation part}
\textbf{Implement:} a function \lstinline|gfp_poly_divide(dividend, divisor, p)| that performs polynomial long division over $GF(p)$. Both polynomials are represented as lists of integer coefficients from highest degree to lowest (e.g., $2x^2 + 1 \rightarrow$ \lstinline|[2, 0, 1]|). You will need to calculate the modular inverse of the divisor's leading coefficient.

\begin{center}
\renewcommand{\arraystretch}{2.0}
\begin{tabular}{r | c | c | c | c}
\toprule
\textbf{Math Symbol} & $f(x)$ & $g(x)$ & $p$ & $(q(x), r(x))$ \\ \midrule
\textbf{Argument} & \lstinline|dividend| & \lstinline|divisor| & \lstinline|p| & \lstinline|return| \\ \midrule
\textbf{Type} & \lstinline|list[int]| & \lstinline|list[int]| & \lstinline|int| & \lstinline|tuple(list[int], list[int])| \\
\bottomrule
\end{tabular}
\end{center}

\section*{Problem 2: Primitive Polynomials}

\subsubsection*{Mathematical part}
Reed--Solomon codes require a field whose nonzero elements can be generated by repeated powers of a single element $\alpha^k$. A primitive polynomial is precisely an irreducible polynomial whose root acts as such a generator for the multiplicative group of order $p^n - 1$.

\textbf{Worked Example: Proving Primitivity} \\
Is $P(x) = x^2 + x + 1$ primitive over $GF(2)$? Here, $p=2, n=2$. The group order is $2^2 - 1 = 3$. We must check if $x^3 \equiv 1 \pmod{P(x)}$.
\begin{enumerate}
    \item We know $x^2 + x + 1 \equiv 0 \implies x^2 \equiv x + 1 \pmod 2$.
    \item Calculate $x^3 = x \cdot x^2 = x(x + 1) = x^2 + x$.
    \item Substitute $x^2$ again: $(x + 1) + x = 2x + 1$.
    \item Modulo 2, $2x = 0$, leaving just $1$.
\end{enumerate}
Since $x^3 \equiv 1 \pmod{x^2+x+1}$, it is primitive.

\subsubsection*{Implementation part}
\textbf{Implement:} a function \lstinline|is_primitive(poly, p, n)| that returns \lstinline|True| if the given polynomial is primitive. You must verify that $x^{(p^n - 1)} \equiv 1 \pmod{\text{poly}}$, and $x^k \not\equiv 1 \pmod{\text{poly}}$ for all proper divisors $k$ of $p^n - 1$.

\begin{center}
\renewcommand{\arraystretch}{2.0}
\begin{tabular}{r | c | c | c | c}
\toprule
\textbf{Math Symbol} & $P(x)$ & $p$ & $n$ & Primitive? \\ \midrule
\textbf{Argument} & \lstinline|poly| & \lstinline|p| & \lstinline|n| & \lstinline|return| \\ \midrule
\textbf{Type} & \lstinline|list[int]| & \lstinline|int| & \lstinline|int| & \lstinline|bool| \\
\bottomrule
\end{tabular}
\end{center}





\section*{Problem 3: Extension Fields GF(p\textsuperscript{n}) Tables}

\subsubsection*{Mathematical part}
An extension field is created by adjoining a root of an irreducible polynomial. To multiply quickly in $GF(p^n)$, we map every polynomial combination of $\alpha$ to an integer, building Logarithm and Exponential lookup tables.

\textbf{Integer Representation:} \
A polynomial in $\alpha$ over $GF(2)$ is mapped to an integer by treating its coefficients as the bits of a binary number. For example, the polynomial $1\cdot\alpha^2 + 0\cdot\alpha^1 + 1\cdot\alpha^0$ translates to the binary sequence $101_2$, which equals the decimal integer $5$.

\textbf{Worked Example: Generating Tables for $GF(2^3)$} \
Using the primitive polynomial $m(x)=x^3+x+1$, we introduce a symbol $\alpha$ which is defined as a root of this polynomial. Therefore:


$$\alpha^3 + \alpha + 1 = 0$$


Because we are working in $GF(2)$, addition and subtraction are the exact same operation (since $1 + 1 = 0$, meaning $1 = -1$). We can move terms across the equals sign without changing their signs:


$$\alpha^3 = \alpha + 1$$


Now we can generate the entire field by multiplying by $\alpha$ and reducing any $\alpha^3$ terms we encounter:
\begin{itemize}
\item $\alpha^0 = 1$ \hfill (Binary: $001_2 \implies$ Integer: 1)
\item $\alpha^1 = \alpha$ \hfill (Binary: $010_2 \implies$ Integer: 2)
\item $\alpha^2 = \alpha^2$ \hfill (Binary: $100_2 \implies$ Integer: 4)
\item $\alpha^3 = \alpha + 1$ \hfill (Binary: $011_2 \implies$ Integer: 3)
\item $\alpha^4 = \alpha(\alpha + 1) = \alpha^2 + \alpha$ \hfill (Binary: $110_2 \implies$ Integer: 6)
\end{itemize}
From this iterative sequence, your code will populate the tables (e.g., \lstinline|exp_table[4] = 6| and \lstinline|log_table[6] = 4|).

\subsubsection*{Implementation part}
\textbf{Implement:} a function \lstinline|generate_gfpn_tables(p, n, primitive_poly)| that constructs the Exponential and Logarithmic lookup tables iteratively by multiplying by $x$ and finding the remainder modulo the primitive polynomial.

\begin{center}
\renewcommand{\arraystretch}{2.0}
\begin{tabular}{r | c | c | c | c}
\toprule
\textbf{Math Symbol} & $p$ & $n$ & $P(x)$ & (Exp Table, Log Table)\\ \midrule
\textbf{Argument} & \lstinline|p| & \lstinline|n| & \lstinline|primitive_poly| & \lstinline|return| \\ \midrule
\textbf{Type} & \lstinline|int| & \lstinline|int| & \lstinline|list[int]| & \lstinline|tuple(list[int], list[int])| \\
\bottomrule
\end{tabular}
\end{center}





\section*{Problem 4: Polynomial Arithmetic in GF(p\textsuperscript{n})}

\subsubsection*{Mathematical part}
The ring $GF(p^n)[x]$ behaves similarly to ordinary polynomial algebra, but the coefficients themselves are elements of the extension field. You must use your Log/Exp tables to multiply or divide coefficients.

\textbf{Worked Example: Polynomial Multiplication in $GF(2^3)$} \\
Let us multiply $A(x) = \alpha x + 1$ and $B(x) = x + \alpha^2$.
\begin{enumerate}
    \item Expand: $(\alpha x)(x) + (\alpha x)(\alpha^2) + (1)(x) + (1)(\alpha^2)$.
    \item Simplify coefficients: $\alpha x^2 + \alpha^3 x + 1 x + \alpha^2$.
    \item Group terms: $\alpha x^2 + (\alpha^3 + 1)x + \alpha^2$.
    \item Since $\alpha^3 = \alpha + 1$ (from Problem 3), we know $\alpha^3 + 1 = \alpha$.
\end{enumerate}
\textbf{Result:} $\alpha x^2 + \alpha x + \alpha^2$.

\subsubsection*{Implementation part}
\textbf{Implement:} functions \lstinline|gfpn_poly_multiply| and \lstinline|gfpn_poly_remainder| utilizing the lookup tables you generated in Problem 3.

\begin{center}
\renewcommand{\arraystretch}{2.0}
\begin{tabular}{r | c | c | c}
\toprule
\textbf{Math Symbol} & $f(x)$ & $g(x)$ & $f(x) \pmod{g(x)}$ \\ \midrule
\textbf{Argument} & \lstinline|dividend| & \lstinline|divisor| & \lstinline|return| \\ \midrule
\textbf{Type} & \lstinline|list[int]| & \lstinline|list[int]| & \lstinline|list[int]| \\
\bottomrule
\end{tabular}
\end{center}

\section*{Problem 5: Reed-Solomon Encoding}

\subsubsection*{Mathematical part}
Reed--Solomon codes are evaluation codes. The generator polynomial $g(x)=\prod_{i=0}^{2t-1}(x-\alpha^i)$ forces every valid codeword to vanish at predetermined field elements. If a code has $2t$ parity symbols, then it can correct up to $t$ symbol errors.

\textbf{Worked Example: Building the Generator} \\
To correct 1 error ($t=1$), we need $2t = 2$ roots: $\alpha^0$ and $\alpha^1$.
$g(x) = (x - \alpha^0)(x - \alpha^1) = (x - 1)(x - \alpha)$. \\
Expanding over $GF(2^3)$: $x^2 - \alpha x - 1x + \alpha$. \\
In $GF(2)$, subtraction is addition, and $1 + \alpha = \alpha^3$. \\
\textbf{Result:} $g(x) = x^2 + \alpha^3 x + \alpha$.

\subsubsection*{Implementation part}
\textbf{Implement:} a function \lstinline|get_generator_poly(num_ec_bytes, log_table, exp_table, p, n)| that calculates the roots and expands the polynomial.

\begin{center}
\renewcommand{\arraystretch}{2.0}
\begin{tabular}{r | c | c }
\toprule
\textbf{Math Symbol} & $2t$ & $g(x)$ \\ \midrule
\textbf{Argument} & \lstinline|num_ec_bytes| & \lstinline|return| \\ \midrule
\textbf{Type} & \lstinline|int| & \lstinline|list[int]| \\
\bottomrule
\end{tabular}
\end{center}




\section*{Problem 6: Decoding I -- Syndromes}

\subsubsection*{Mathematical part}
When a message is received, it may contain errors. We represent the received message as a polynomial $r(x) = c(x) + e(x)$, where $c(x)$ is the original valid codeword and $e(x)$ represents the errors. Because a valid codeword evaluates to zero at the roots of the generator polynomial ($\alpha^0, \alpha^1, \dots, \alpha^{2t-1}$), any non-zero result when evaluating $r(x)$ at these roots must be caused purely by the errors. These results are called \textbf{syndromes}.

\textbf{Worked Example: Calculating Syndromes} \\
Assume we are working in $GF(2^3)$ and received $r(x) = x^2 + \alpha x + 1$. We need to correct 1 error ($t=1$), so we calculate $2t=2$ syndromes at roots $\alpha^0$ and $\alpha^1$.
\begin{enumerate}
    \item \textbf{Evaluate at $\alpha^0 = 1$:} \\
    $S_1 = r(1) = (1)^2 + \alpha(1) + 1 = 1 + \alpha + 1 = \alpha$.
    \item \textbf{Evaluate at $\alpha^1 = \alpha$:} \\
    $S_2 = r(\alpha) = (\alpha)^2 + \alpha(\alpha) + 1 = \alpha^2 + \alpha^2 + 1 = 1$.
\end{enumerate}
\textbf{Result:} The syndrome sequence is $[\alpha, 1]$. Because they are not zero, an error exists.

\subsubsection*{Implementation part}
\textbf{Implement:} \lstinline|calculate_syndromes(message, num_ec, log_table, exp_table, p, n)|. You must evaluate the received message polynomial at each root $\alpha^0, \dots, \alpha^{\text{num\_ec}-1}$.

\begin{center}
\renewcommand{\arraystretch}{2.0}
\begin{tabular}{r | c | c | c }
\toprule
\textbf{Math Symbol} & $r(x)$ & $2t$ & $[S_1, S_2, \dots, S_{2t}]$ \\ \midrule
\textbf{Argument} & \lstinline|message| & \lstinline|num_ec| & \lstinline|return| \\ \midrule
\textbf{Type} & \lstinline|list[int]| & \lstinline|int| & \lstinline|list[int]| \\
\bottomrule
\end{tabular}
\end{center}



\section*{Problem 7: Decoding II -- Berlekamp-Massey}

\subsubsection*{Mathematical part}
The Berlekamp-Massey algorithm finds the shortest Linear Feedback Shift Register (LFSR) that generates your syndrome sequence. It is an iterative state machine. You must maintain the following state variables:
\begin{itemize}
    \item $C(x)$: The current Error Locator Polynomial (starts as $1$).
    \item $B(x)$: A copy of the last $C(x)$ before the assumed number of errors ($L$) was updated (starts as $1$).
    \item $L$: The current assumed number of errors (starts at $0$).
    \item $m$: The shift counter. The number of iterations since $L$ was last updated (starts at $1$).
    \item $b$: The discrepancy from the last time $L$ was updated (starts at $1$).
\end{itemize}

\textbf{The Iterative Algorithm:} \\
For each syndrome $S_i$ (where $i$ goes from $0$ to $2t-1$):
\begin{enumerate}
    \item Calculate the current discrepancy $d$:
    $$d = S_i + \sum_{j=1}^{L} C_j S_{i-j}$$
    \item \textbf{If $d == 0$:} The current $C(x)$ is still valid. Just increment the shift counter: $m = m + 1$.
    \item \textbf{If $d \neq 0$:} We must adjust $C(x)$.
    \begin{itemize}
        \item Calculate the adjusted polynomial: $T(x) = C(x) - (d \cdot b^{-1}) \cdot x^m \cdot B(x)$
        \item \textbf{If $2L \le i$:} We need to increase our assumed number of errors.
            \begin{itemize}
                \item $L = i + 1 - L$
                \item $B(x) = C(x)$ \quad \textit{(Save the old C(x) before updating it)}
                \item $b = d$
                \item $m = 1$
            \end{itemize}
        \item \textbf{Else:} We don't update $L$, $B(x)$, or $b$. Just increment $m = m + 1$.
        \item Finally, apply the adjustment: $C(x) = T(x)$.
    \end{itemize}
\end{enumerate}

\textbf{Worked Example: Tracing the State Machine} \\
Let $S = [\alpha, 1]$. Initial state: $C=[1], B=[1], L=0, m=1, b=1$.
\begin{itemize}
    \item \textbf{Loop $i=0$ ($S_0 = \alpha$):}
    \begin{itemize}
        \item $d = S_0 = \alpha$. ($d \neq 0$).
        \item $T(x) = [1] - (\alpha \cdot 1^{-1}) \cdot x^1 \cdot [1] = 1 - \alpha x$.
        \item $2(0) \le 0$ is True. Update: $L = 0 + 1 - 0 = 1$. $B(x) = 1$. $b = \alpha$. $m = 1$.
        \item $C(x) = 1 - \alpha x$. (In $GF(2)$, this is $1 + \alpha x$).
    \end{itemize}
    \item \textbf{Loop $i=1$ ($S_1 = 1$):}
    \begin{itemize}
        \item $d = S_1 + C_1 S_0 = 1 + (\alpha)(\alpha) = 1 + \alpha^2$. ($d \neq 0$).
        \item $T(x) = (1 + \alpha x) - ((1+\alpha^2) \cdot \alpha^{-1}) \cdot x^1 \cdot [1]$.
        \item $2(1) \le 1$ is False. Update: $m = 2$.
        \item $C(x) = T(x)$.
    \end{itemize}
\end{itemize}

\subsubsection*{Implementation part}
\textbf{Implement:} \lstinline|berlekamp_massey(syndromes, log_table, exp_table, p, n)|. Use your $GF(p^n)$ multiplication, division, and addition functions to execute the state machine described above. Pay special attention to list padding when subtracting the shifted $B(x)$ from $C(x)$.

\begin{center}
\renewcommand{\arraystretch}{2.0}
\begin{tabular}{r | c | c }
\toprule
\textbf{Math Symbol} & $[S_0, S_1, \dots, S_{2t-1}]$ & $C(x)$ \\ \midrule
\textbf{Argument} & \lstinline|syndromes| & \lstinline|return| \\ \midrule
\textbf{Type} & \lstinline|list[int]| & \lstinline|list[int]| \\
\bottomrule
\end{tabular}
\end{center}



\section*{Problem 8: Decoding III -- Chien Search \& Forney}

\subsubsection*{Mathematical part}
Now that we have the Error Locator Polynomial $\Lambda(x)$, we must find its roots (Chien Search) and calculate how severely the bits were corrupted (Forney Algorithm).

\textbf{Step 1: Chien Search} \\
This is a brute-force evaluation. We plug every possible inverse field element $\alpha^{-i}$ into $\Lambda(x)$. If the result is $0$, there is an error at index $i$ in the message.
\textit{Example:} If $\Lambda(x) = 1 + \alpha^3 x$, testing $x = \alpha^{-3}$ gives $1 + \alpha^3(\alpha^{-3}) = 1 + 1 = 0$. The root is $\alpha^{-3}$, meaning an error exists at position $3$.

\textbf{Step 2: Forney Algorithm} \\
To find the exact error magnitude, we need the Error Evaluator Polynomial:
$$\Omega(x) \equiv S(x)\Lambda(x) \pmod{x^{2t}}$$
The magnitude of the error at position $i$ (corresponding to root $X_i^{-1}$) is given by the formal derivative of the locator polynomial $\Lambda'(x)$:
$$e_i = - \frac{X_i \Omega(X_i^{-1})}{\Lambda'(X_i^{-1})}$$

\subsubsection*{Implementation part}
\textbf{Implement:} \lstinline|chien_search| to find the integer indices of the errors, and \lstinline|forney_algorithm| to calculate a dictionary mapping those indices to their exact error magnitudes.

\begin{center}
\renewcommand{\arraystretch}{2.0}
\begin{tabular}{r | c | c | c | c}
\toprule
\textbf{Math Symbol} & $\Lambda(x)$ & message length & Error Positions \\ \midrule
\textbf{Argument} & \lstinline|err_loc| & \lstinline|msg_len| & \lstinline|return| \\ \midrule
\textbf{Type} & \lstinline|list[int]| & \lstinline|int| & \lstinline|list[int]| \\
\bottomrule
\end{tabular}
\end{center}

\vspace{0.5cm}

\begin{center}
\renewcommand{\arraystretch}{2.0}
\begin{tabular}{r | c | c | c | c}
\toprule
\textbf{Math Symbol} & $[S_1, \dots]$ & $\Lambda(x)$ & Error Positions & Magnitudes \\ \midrule
\textbf{Argument} & \lstinline|syndromes| & \lstinline|err_loc| & \lstinline|err_pos| & \lstinline|return| \\ \midrule
\textbf{Type} & \lstinline|list[int]| & \lstinline|list[int]| & \lstinline|list[int]| & \lstinline|dict{pos: mag}| \\
\bottomrule
\end{tabular}
\end{center}




\section*{Problem 9: The QR Matrix}

\subsubsection*{Mathematical part}
The QR matrix is simply the final representation of a Reed--Solomon encoded bitstream. To ensure cameras can read it regardless of lighting, the ISO specification applies a data mask.

\textbf{Worked Example: Masking} \\
QR codes use XOR masks to break up large blocks of identical colors. If Mask 0 is defined as $(row + col) \bmod 2 == 0$, and you are about to place a `1` (black) bit at coordinate $(row=8, col=6)$:
\begin{enumerate}
    \item Calculate $(8 + 6) \bmod 2 = 14 \bmod 2 = 0$.
    \item The mask condition is True.
    \item XOR the bit: $1 \oplus 1 = 0$. The pixel becomes white.
\end{enumerate}

\subsubsection*{Implementation part}
\textbf{Implement:} \lstinline|build_qr_matrix(version, all_bits, get_format_string)| to route your robustly encoded Reed-Solomon bitstream into the physical 2D matrix in a zigzag pattern, applying Mask 0. Use the provided \lstinline|draw_finder| logic to avoid corrupting the tracking squares.

\textbf{Test:} Run \lstinline|qr_code_demo.py| to test your final matrix with a real smartphone camera!

\end{document}
